\documentclass{report}
\usepackage[letterpaper,portrait,margin=2cm]{geometry}
\usepackage{amsmath,amssymb,amsthm,mathtools,dsfont,xfrac,gensymb}
\usepackage{enumitem}
\usepackage[french]{babel}
\usepackage{titlesec}
\usepackage{fancyhdr}
\usepackage{hyperref}
\usepackage{tikz}
\usetikzlibrary{calc, positioning, 3d}

\hypersetup{
	colorlinks=false
}

\allowdisplaybreaks

\fancyfoot{}
\fancyhead{}
\fancyhead[l]{MAT603 - G\'eom\'etrie diff\'erentielle}
\fancyhead[r]{\leftmark}
\fancyfoot[c]{\thepage}
\renewcommand{\headrulewidth}{.4pt}
\renewcommand{\footrulewidth}{0pt}

\setcounter{tocdepth}{3}

\pagestyle{fancy}

\title{MAT603 - G\'eom\'etrie diff\'erentielle

	Donn\'e par Maxence Mayrand}
\author{Julien Houle}
\date{Automne 2026}

\renewcommand{\thesection}{ \arabic{chapter}.\arabic{section}}
\renewcommand{\thesubsection}{}
\renewcommand{\thesubsubsection}{}

\titleformat{\chapter}[hang]{\bfseries\huge\centering}{Chapitre \arabic{chapter}}{1em}{}[]
\titleformat{\section}[hang]{\bfseries\large}{Section\thesection}{1em}{}[]
\titleformat{\subsection}[hang]{\bfseries\normalsize}{}{0pt}{}[]
\titleformat{\subsubsection}[hang]{\slshape\normalsize}{}{0pt}{}[]

\newtheorem*{thm}{Th\'eor\`eme}
\newtheorem*{lem}{Lemme}
\newtheorem*{prop}{Proposition}
\newtheorem*{coro}{Corollaire}
\theoremstyle{definition}
\newtheorem*{dfn}{D\'efinition}
\theoremstyle{remark}
\newtheorem*{exem}{Exemple}
\newtheorem*{exer}{Exercice}
\newtheorem*{nota}{Notation}
\newtheorem*{rmq}{Remarque}
\newtheorem*{rapp}{Rappel}

\setlist{nosep}

\newenvironment*{subproof}[1][D\'emonstration]
{
\begin{proof}[#1]

}
{
\renewcommand{\qedsymbol}{$\blacksquare$}
\end{proof}
\renewcommand{\qedsymbol}{$\square$}
}

\def\lte{\leqslant}
\def\gte{\geqslant}

\def\NN{\mathbb{N}}
\def\ZZ{\mathbb{Z}}
\def\QQ{\mathbb{Q}}
\def\RR{\mathbb{R}}
\def\CC{\mathbb{C}}

\def\id{\mathds{1}}

\def\sech{\operatorname{sech}}
\def\csch{\operatorname{csch}}
\def\coth{\operatorname{coth}}
\def\arcsinh{\operatorname{arcsinh}}
\def\arccosh{\operatorname{arccosh}}
\def\arctanh{\operatorname{arctanh}}
\def\arcsech{\operatorname{arcsech}}
\def\arccsch{\operatorname{arccsch}}
\def\arccoth{\operatorname{arccoth}}

\newcommand*{\abs}[1]{\left|#1\right|}
\newcommand*{\norme}[1]{\left\|#1\right\|}

\begin{document}
	\maketitle
	\pagenumbering{roman}
	\tableofcontents
	\newpage
	\pagenumbering{arabic}
	\phantomsection\addcontentsline{toc}{chapter}{Survol du cours}
	\chapter*{Survol du cours}
	\phantomsection\addcontentsline{toc}{section}{Courbes}
	\section*{Courbes dans $\RR^3$ ou $\RR^2$}
	La courbure est caract\'eris\'ee par $\kappa=\sfrac{1}{r}$. La torsion est caract\'eris\'ee par $\tau=\vec N \cdot \vec T$.

	\begin{thm}[Th\'eor\`eme fondamental des courbes]
		La courbure et la torsion d\'eterminent une courbe compl\`etement, \`a transformation rigide pr\`es.
	\end{thm}

	\phantomsection\addcontentsline{toc}{section}{Surfaces}
	\section*{Surfaces dans $\RR^3$}
	Une cr\'eature bidimensionnelle peut ``vivre'' sur une surface, se d\'eplacer dessus et mesurer des distances et des angles.

	La g\'eom\'etrie intrins\`eque est d\'ecrite par un produit scalaire sur des vecteurs tangents \`a la surface.

	On peut intersecter la surface avec des plans et mesurer la courbure de la courbe \`a l'intersection du plan et de la surface.

	$k_1$ est le max des courbures et $k_2$ est le min des courbures. $K=k_1k_2$ est la courbure de Gauss.

	\begin{thm}[Theorema egregium (Gauss)]
		$K$ est intrins\`eque.
	\end{thm}

	\begin{thm}[Gauss-Bonnet]
		$\alpha+\beta+\gamma = \pi + \int_{\triangle}K$.
	\end{thm}

	Une cr\'eature se d\'eplace en ligne ``droite'' sur la surface en laissant une corde le long du chemin. Quand elle revient au point de d\'epart, elle tire sur les deux c\^ot\'es de la corde. Si la corde revient, le chemin ne passe pas par un trou. Si la corde est prise, le chemin passe par un trou.

	\chapter{Courbes}
	Soit $I \subseteq \RR$ un intervalle.

	Une application $f:I \to \RR^3$ d\'efinie par $f(t) = (f_1(t), f_2(t), f_3(t))$ est \emph{diff\'erentiable} si $f_1',f_2',f_3'$ existent sur $I$. $f'(t) = (f_1'(t), f_2'(t), f_3'(t))$. $F$ est $C^k$ si $f', f'', \cdots, f^{(k)}$ existent et sont continues. $F$ est \emph{lisse} si $f$ est $C^k$, $\forall k\in \NN$.

	Une \emph{courbe param\'etr\'ee} est une application $\vec\alpha:I \to \RR^3$, o\`u $I$ est un intervalle et $\vec\alpha$ est $C^3$. $\vec\alpha$ est \emph{r\'eguli\`ere} si $\vec\alpha'(t)\neq0, \forall t\in I$.

	\begin{exem}
		\begin{enumerate}
			\item Droite: $\vec P,\vec Q \in \RR^3$, $\vec\alpha(t) = \vec P + t(\vec Q-\vec P)$, $I=\RR$;
			\item Cercle: $\vec\alpha(t) = (a\cos t,a\sin t)$, $I=[0,2\pi]$;
			\item Ellipse: $\left\{ (x,y)\in\RR^2 ~\middle|~ \frac{x^2}{a^2}+\frac{y^2}{b^2}=1 \right\}$. $T : \RR^2\to\RR^2 : (x,y)\mapsto(ax,by)$ envoie le cercle \`a l'ellipse. $\vec\alpha(t) = T(\cos t,\sin t) = (a\cos t,b\sin t)$, $I=[0,2\pi]$;
			\item Cyclo\"ide: La trajectoire suivie par un point sur un cercle de rayon $a$ lorsque le cercle roule sur une droite.

			On a
			\begin{itemize}
				\item $l=at$,
				\item le centre du cercle est \`a $(l,a)=(at,a)$,
				\item $\vec\alpha(t) = (at,a)+(-a\sin t,-a\cos t) = (a(t-\sin t), a(1-\cos t))$.
			\end{itemize}
			\item H\'elice: $\vec\alpha(t) = (a\cos t, a\sin t, bt)$, $a,b>0$.
		\end{enumerate}
	\end{exem}

	\begin{rapp}
		Fonctions hyperboliques.
		\begin{align*}
			\begin{split}
				\cosh t&= \dfrac{e^t+e^{-t}}{2}
			\end{split}
			&
			\begin{split}
				\sinh t&= \dfrac{e^t-e^{-t}}{2}
			\end{split}
			\\[2ex]
			\begin{split}
				\tanh t&= \dfrac{\sinh t}{\cosh t}
			\end{split}
			&
			\begin{split}
				\sech t&= \dfrac{1}{\cosh t}
			\end{split}
		\end{align*}
	\end{rapp}

	\begin{enumerate}[start=6]
		\item Cha\^inette: $\vec\alpha(t)=(t, c\cosh(\frac{t}{c}))$, o\`u $c>0$;
		\item Tractrice: Points \`a distance constante d'un point sur une droite en direction d'un autre point. $\vec\alpha'(t)$ est dirig\'e du point externe vers le point associ\'e \`a $t$ sur la droite.

		Soit $\theta(t)$ l'angle entre la droite et la direction du point externe.

		$\vec\alpha'(t) = c (\cos\theta(t), \sin\theta(t))$, $\vec\alpha(t) = (x(t),y(t)) = (t,0) + (\cos\theta(t), \sin\theta(t))$, $x(t) = t+\cos\theta(t)$, $y(t) = \sin\theta(t)$.

		$(x'(t),y'(t)) = c (\cos\theta(t), \sin\theta(t))$, donc $(1-\sin\theta(t)\theta'(t), \cos\theta(t)\theta'(t)) = c (\cos\theta(t), \sin\theta(t))$.

		\begin{align*}
			\begin{split}
				\dfrac{y'(t)}{x'(t)}&= \dfrac{c\sin\theta(t)}{c\cos\theta(t)}\\
				&= \tan\theta(t)
			\end{split}
			&
			\begin{split}
				\dfrac{y'(t)}{x'(t)}&= \dfrac{\cos\theta\theta'}{1-\sin\theta\theta'}
			\end{split}
		\end{align*}
		donc
		\begin{align*}
			\cos\theta\theta'&= \dfrac{\sin\theta}{\cos\theta} - \dfrac{\sin^2\theta}{\cos\theta}\theta'\\
			\cos^2\theta\theta'&= \sin\theta-\sin^2\theta\theta'\\
			(\cos^2\theta+\sin^2\theta)\theta'&= \sin\theta\\
			\theta'(t)&= \sin\theta(t)
		\end{align*}

		\begin{align*}
			\int \dfrac{d\theta}{\sin\theta}&= \int dt
		\end{align*}
		or
		\begin{align*}
			\int \csc\theta d\theta&= \int \csc\theta \dfrac{\csc\theta+\cot\theta}{\csc\theta+\cot\theta} d\theta\\
			&= \int \dfrac{\csc^2\theta+\csc\theta\cot\theta}{\csc\theta+\cot\theta} d\theta\\
			&= -\int \dfrac{d(\csc\theta+\cot\theta)}{\csc\theta+\cot\theta}\\
			&= -\ln(\csc\theta+\cot\theta)+C
		\end{align*}
		donc
		\begin{align*}
			t&= -\ln(\csc\theta+\cot\theta)+C\\
			\shortintertext{comme $\theta(0)=\sfrac{\pi}{2}$,}
			&= -\ln(\csc\theta+\cot\theta)\\
			&= -\ln\left(\cot\frac{\theta}{2}\right)\\
			&= \ln\left(\tan\frac{\theta}{2}\right)
		\end{align*}

		Donc, $\vec\alpha(t) = (t+\cos\theta(t), \sin\theta(t))$, alors $\vec\beta(\theta) = (\ln\tan\frac{\theta}{2} + \cos\theta, \sin\theta)$ est une param\'etrisation en fonction de $\theta$.

		\begin{align*}
			t&= \ln\tan\frac{\theta}{2}\\
			e^t&= \tan\frac{\theta}{2}\\
			\sin\theta&= 2\sin\frac{\theta}{2}\cos\frac{\theta}{2}\\
			&= 2\tan\frac{\theta}{2}\cos^2\frac{\theta}{2}\\
			&= \dfrac{2\tan\frac{\theta}{2}}{\frac{1}{\cos^2\frac{\theta}{2}}}\\
			&= \dfrac{2\tan\frac{\theta}{2}}{1+\tan^2\frac{\theta}{2}}\\
			&= \dfrac{2e^t}{1+e^{2t}}\\
			&= \dfrac{2}{e^{-t}+e^t}\\
			&= \sech t
		\end{align*}
		\begin{align*}
			1+\cos\theta&= 2\cos^2\frac{\theta}{2}\\
			\cos\theta&= 2\cos^2\frac{\theta}{2}-1\\
			&= \cos^2\frac{\theta}{2}-\sin^2\frac{\theta}{2}\\
			&= \dfrac{1-\tan^2\frac{\theta}{2}}{\frac{1}{\cos^2\frac{\theta}{2}}}\\
			&= \dfrac{1-\tan^2\frac{\theta}{2}}{1+\tan^2\frac{\theta}{2}}\\
			&= \dfrac{1-e^{2t}}{1+e^{2t}}\\
			&= \tanh t
		\end{align*}

		Ainsi, $\vec\alpha(t) = (t-\tanh t, \sech t)$.
	\end{enumerate}

	\begin{dfn}
		Soit $\vec\alpha:[a,b]\to\RR^3$ une courbe param\'etr\'ee. La \emph{longueur d'arc} de $a$ \`a $t\in[a,b]$ est $s(t) = \int_{a}^{t}\|\vec\alpha'(u)\|du$.

		Soit $P = \{t_0, t_1, \cdots, t_n\}$ une partition de $[a,b]$.

		$l(\vec\alpha, P) = \sum_{i=1}^{n} \norme{\vec\alpha(t_i) - \vec\alpha(t_{i-1})}$ est une approximation de la longueur de $\vec\alpha$. Plus $P$ est fine, plus l'approximation est bonne. La longueur de $\vec\alpha$ est $\sup\limits_P \sum_{i=1}^{n} \norme{\vec\alpha(t_i) - \vec\alpha(t_{i-1})}$. $\vec\alpha'(t_i) = \lim\limits_{t\to t_i} \frac{\vec\alpha(t) - \vec\alpha(t_i)}{t-t_i} \approx \frac{\vec\alpha(t_{i+1}) - \vec\alpha(t_i)}{t_{i+1}-t_i}$. La longueur de $\vec\alpha$ est approximativement $\sup\limits_P \sum_{i=1}^{n} \norme{\vec\alpha(t_{i-1})} (t_i-t_{i-1}) = \sup\limits_P \sum_{i=1}^{n} \norme{\vec\alpha'(t_{i-1})} (t_i-t_{i-1})$. Alors, la longueur de $\vec\alpha$ est $\int_{a}^{b} \norme{\vec\alpha'(u)} du$.

		On dit que $\vec\alpha$ est \emph{param\'etr\'ee par la longueur d'arc} si $s(t)=t$.
	\end{dfn}

	\begin{rmq}
		$s(t)=t \Leftrightarrow a=0$ et $\norme{\vec\alpha'(t)}=1, \forall t\in[a,b]$.
		\begin{proof}
			\begin{subproof}[$\Rightarrow$]
				$t=s(t)=\int_{a}^{t}\norme{\alpha'(u)}du$. En d\'erivant de chaque c\^ot\'e, on a $1=\norme{\alpha(t)}$. En \'evaluant en $t=a$, on a $a=\int_{a}^{a}\norme{\alpha'(u)}du=0$.
			\end{subproof}
			\begin{subproof}[$\Leftarrow$]
				$s(t) = \int_{0}^{t} \norme{\alpha'(u)}du = \int_{0}^{t}1du = t$.
			\end{subproof}
		\end{proof}
	\end{rmq}

	\begin{exem}
		$\vec\alpha(t) = (r\cos t, r\sin t)$, avec $t\in[0,2\pi]$.
		\begin{align*}
			\norme{\vec\alpha'(t)}&= \norme{(-r\sin t, r\cos t)}\\
			&= \sqrt{r^2\sin^2t+r^2\cos^2t}\\
			&= r
		\end{align*}
		\begin{align*}
			s(t)&= \int_{0}^{t}\norme{\vec\alpha'(u)}du\\
			&= \int_{0}^{t}rdu\\
			&= rt
		\end{align*}

		Posons $\vec\beta(t) = \vec\alpha(\sfrac{t}{r}) = (r\cos\frac{t}{r}, r\sin\frac{t}{r})$. $\norme{\vec\beta'(t) = \sqrt{\cos^2\frac{t}{r} + \sin^2\frac{t}{r}}} = 1$. $\vec\beta(s) = (r\cos\frac{s}{r}, r\sin\frac{s}{r})$ est une param\'etrisation par longueur d'arc.
	\end{exem}

	Toute courbe r\'eguli\`ere peut \^etre reparam\'etr\'ee par longueur d'arc.
	\begin{proof}[En effet]
		$s(t) = \int_{a}^{t} \norme{\vec\alpha'(u)} du$ est strictement croissante, car $\norme{\vec\alpha'(u)}>0, \forall u$.

		Donc, il existe un inverse $f:[0,L]\to[a,b]$, c'est-\`a-dire, $s(f(t))=t$ et $f(s(t))=t$.

		Soit $\vec\beta(t)=\vec\alpha(f(t))$. $\vec\beta'(t) = \vec\alpha'(f(t)) f'(t)$, donc $\vec\beta(u) = \vec\alpha'(f(u))f'(u) = \frac{\vec\alpha'(f(u))}{s'(f(u))} = \frac{\vec\alpha'(f(u))}{\norme{\vec\alpha'(f(u))}}$. Ainsi, $\norme{\vec\beta'(u)}=1$ et $\vec\beta$ est param\'etr\'ee par longueur d'arc.
	\end{proof}

	\begin{exem}
		$\vec\alpha(t) = (t, \cosh t, \sinh t)$, avec $t\in[0,1]$. Reparam\'etrisons par longueur d'arc.
		\begin{align*}
			s(t)&= \int_{0}^{t} \norme{\vec\alpha'(u)}du\\
			&= \int_{0}^{t} \norme{(1, \sinh u, \cosh u)} du\\
			&= \int_{0}^{t} \sqrt{1+\sinh^2u+\cosh^2u} du\\
			&= \int_{0}^{t} \sqrt{2\cosh^2u} du\\
			&= \sqrt{2} \int_{0}^{t} \cosh u du\\
			&= \left. \sqrt{2} \sinh u \right|_{0}^{t}\\
			s(t)&= \sqrt{2}\sinh t
		\end{align*}
		donc
		\begin{align*}
			\frac{s}{\sqrt{2}}&= \sinh t\\
			\arcsinh\frac{s}{\sqrt{2}}&= t
		\end{align*}

		\begin{align*}
			\vec\beta(s)&= \vec\alpha(\arcsinh\frac{s}{\sqrt{2}})\\
			&= \left( \arcsinh\frac{s}{\sqrt{2}}, \cosh\left( \arcsinh\frac{s}{\sqrt{2}} \right), \frac{s}{\sqrt{2}} \right)\\
			&= \left( \arcsinh\frac{s}{\sqrt{2}}, \sqrt{1+\frac{s^2}{2}}, \frac{s}{\sqrt{2}} \right)
		\end{align*}
	\end{exem}

	\section{Rep\`ere de Frenet}
	\begin{lem}
		Soient $f,g: ]a,b[ \to \RR^3$ deux courbes diff\'erentiables t.q. $f(t) \cdot g(t)$ est constant $\forall t$, alors $f'(t) \cdot g(t) = -f(t) \cdot g'(t)$.

		En particulier, $\norme{f(t)}$ soit constante est \'equivalent \`a $f(t) \cdot f'(t)=0$, $\forall t$.
		\begin{proof}
			Si $f(t) \cdot g(t)=c$, $\forall t \in ]a,b[$, alors $\frac{d}{dt}(f(t) \cdot g(t)) = 0$, donc $f'(t) \cdot g(t) + f(t) \cdot g'(t)=0$. Ainsi, $f'(t) \cdot g(t) = -f(t) \cdot g'(t)$.

			\begin{align*}
				\norme{f(t)}&= \text{const}\\
				\norme{f(t)}^2&= \text{const}\\
				f(t) \cdot f(t)&= const\\
				f'(t) \cdot f(t)&= -f(t) \cdot f'(t)\\
				2f'(t) \cdot f(t)&= 0\\
				f'(t) \cdot f(t)&= 0
			\end{align*}
		\end{proof}
	\end{lem}

	Soit $\vec\alpha: I \to \RR^3$ une courbe param\'etr\'ee par longueur d'arc.

	$\vec\alpha:I\to\RR^3:s\mapsto\vec\alpha(s)$.

	$\vec T(s) \coloneq \vec\alpha'(s)$ est le \emph{vecteur tangent unitaire}. $\norme{\vec T(s)}=1, \forall s$.

	Par le lemme, $\vec T' \cdot \vec T=0$.

	\begin{dfn}[Courbure]
		On d\'efinit la courbure de $\vec\alpha$ comme $\kappa(s) \coloneq \norme{\vec T'(s)}$.
	\end{dfn}

	\begin{dfn}[Vecteur normal principal]
		Si $\kappa\neq0$, c'est-\`a-dire, $\vec T'\neq0$, on d\'efinit $\vec N(s) = \frac{\vec T'(s)}{\norme{\vec T'(s)}} = \frac{\vec T'(s)}{\kappa(s)}$ le \emph{vecteur normal principal}.
	\end{dfn}

	On a $\vec N \cdot \vec T=0$ et $\norme{\vec N} = \norme{\vec T} = 1$.

	\begin{dfn}[Vecteur binormal]
		$\vec B(s) = \vec T(s) \times \vec N(s)$ est le \emph{vecteur binormal}.
	\end{dfn}

	Le \emph{rep\`ere de Frenet} est la base orthonorm\'ee $\{\vec T(s), \vec N(s), \vec B(s)\}$.

	\begin{dfn}[Torsion]
		La \emph{torsion} est $\tau=\vec N' \cdot \vec B$.
	\end{dfn}

	$\vec T' = \kappa\vec N$.

	$\vec N' = a\vec T+b\vec N+c\vec B$, pour certains $a,b,c$, car $\vec T,\vec N,\vec B$ est une base orthonorm\'ee. $a = \vec N'\cdot\vec T=-\vec N\cdot\vec T' = -\vec N\cdot\kappa\vec N = -\kappa$ par le lemme, car $\vec N\cdot\vec T=0$. $b=\vec N'\cdot\vec N=0$, par le lemme, car $\norme{\vec N}=1$. $c = \vec N'\cdot\vec B = \tau$. Ainsi, $\vec N' = -\kappa\vec T+\tau\vec B$.

	$\vec B' = d\vec T+e\vec N+f\vec B$. $d = \vec B'\cdot\vec T = -\vec B\cdot\vec T' = -\vec B\cdot\kappa\vec N = 0$. $e=\vec B'\cdot\vec N = -\vec B\cdot\vec N' = -\vec B\cdot(-\kappa\vec T+\tau\vec B) = -\tau$. $f = \vec B'\cdot\vec B = 0$, par le lemme. Ainsi, $\vec B'=\tau\vec N$.

	\begin{prop}[Formules de Frenet]
		\begin{align}
			\vec T'&= \kappa\vec N\\
			\vec N'&= -\kappa\vec T+\tau\vec B\\
			\vec B'&= -\tau\vec N
		\end{align}
		o\`u $\kappa=\norme{\vec T'}$ et $\tau=\vec N\cdot\vec B$.
	\end{prop}

	\begin{exem}[H\'elice]
		$\vec\alpha(t) = (a\cos t, a\sin t ,bt)$, $\vec\alpha:[0,\infty[\to\RR^3$.
		\begin{align*}
			s(t)&= \int_{0}^{t} \norme{\vec\alpha'(u)} du\\
			&= \int_{0}^{t} \sqrt{a^2+b^2} du\\
			\shortintertext{en notant $\sqrt{a^2+b^2}=c$,}
			&= tc
		\end{align*}

		$\vec\beta(s) = \vec\alpha(t(s)) = \vec\alpha(\sfrac{s}{c}) = (a\cos\sfrac{s}{c}, a\sin\sfrac{s}{c}, \sfrac{bs}{c})$ est une param\'etrisation par longueur d'arc.
		\begin{align*}
			\begin{split}
				\vec T(s)&= \vec\beta'(s)\\
				&= \left( -\frac{a}{c}\sin\frac{s}{c}, \frac{a}{c}\cos\frac{s}{c}, \frac{b}{c} \right)
			\end{split}
			&
			\begin{split}
				\vec T'(s)&= \left( -\frac{a}{c^2}\cos\frac{s}{c}, -\frac{a}{c^2}\sin\frac{s}{c}, 0 \right)
			\end{split}
			\\[1em]
			\begin{split}
				\kappa(s)&= \norme{\vec T'(s)}\\
				&= \frac{a}{c^2}\\
				&= \frac{a}{a^2+b^2}
			\end{split}
			&
			\begin{split}
				\vec N(s)&= \frac{\vec T'(s)}{\kappa(s)}\\
				&= \left( -\cos\frac{s}{c}, -\sin\frac{s}{c}, 0 \right)
			\end{split}
			\\[1em]
			\begin{split}
				\vec B&= \vec T \times \vec N\\
				&= \left( \frac{b}{c}\sin\frac{s}{c}, -\frac{b}{c}\cos\frac{s}{c}, \frac{a}{c} \right)
			\end{split}
			&
			\begin{split}
				\vec B'&= \left( \frac{b}{c^2}\cos\frac{s}{c}, \frac{b}{c^2}\sin\frac{s}{c}, 0 \right)\\
				&= -\frac{b}{c^2} \left( -\cos\frac{s}{c}, -\sin\frac{s}{c}, 0 \right)\\
				&= -\tau\vec N
			\end{split}
			\\[1em]
			\begin{split}
				\tau&= \frac{b}{c^2}\\
				&= \frac{b}{a^2+b^2}
			\end{split}
			&
			\begin{split}
				b\to0&\Rightarrow\left\{\begin{array}{c}
					\tau\to0\\
					\kappa\to\frac{1}{a}
				\end{array}\right.\\
				b\to\infty&\Rightarrow\left\lbrace\begin{array}{c}
					\tau\to0\\
					\kappa\to0
				\end{array}\right.
			\end{split}
		\end{align*}
	\end{exem}

	Quand $\vec\alpha$ est param\'etr\'ee par longueur d'arc (PPLA),
	\begin{align*}
		\vec T&= \vec\alpha'\\
		\vec N&= \frac{\vec T'}{\kappa}\\
		\vec B&= \vec T \times \vec \vec N
	\end{align*}
	o\`u $\kappa=\norme{\vec T'}$ et $\tau=\vec N' \cdot \vec B$.

	Si $\vec\alpha$ n'est pas PPLA, $s(t) = \int_{a}^{t} \norme{\vec\alpha(u)}du$ est la longueur d'arc.
	\begin{align*}
		\vec T(t)&= \frac{\vec\alpha'(t)}{\norme{\vec\alpha'(t)}}\\
		&= \frac{\vec\alpha'(t)}{v(t)}
	\end{align*}
	o\`u $v(t)=\norme{\vec\alpha'(t)}$ est la vitesse.
	\begin{align*}
		\kappa(t)&= \norme{\frac{d}{dt}\vec T}\\
		&= \norme{\frac{\sfrac{d\vec T}{ds}}{\sfrac{ds}{dt}}}\\
		&= \norme{\frac{\vec T'(t)}{\vec\alpha'(t)}}\\
		&= \frac{\norme{\vec T'(t)}}{v(t)}
	\end{align*}
	\begin{align*}
		\vec N(t)&= \dfrac{\frac{d}{ds}\vec T}{\norme{\frac{d}{ds}\vec T}}\\
		&= \dfrac{\frac{\vec T'}{v(t)}}{\kappa(t)}\\
		&= \frac{\vec T'(t)}{\kappa(t)v(t)}
	\end{align*}
	\begin{align*}
		\vec B(t)&= \vec T(t) \times \vec N(t)
	\end{align*}
	\begin{align*}
		\tau(t)&= \frac{d\vec N}{ds} \cdot \vec B\\
		&= \frac{\vec N'(t)}{v(t)} \cdot \vec B(t)
	\end{align*}

	En g\'en\'eral,
	\begin{align*}
		\vec T&= \frac{\vec\alpha'}{v}\\
		\vec N&= \frac{\vec T'}{\kappa v}\\
		\vec B&= \vec T \times \vec N
	\end{align*}
	o\`u $v=\norme{\vec\alpha'}$, $\kappa=\frac{\norme{\vec T'}}{v}$ et $\tau=\frac{\vec N'\cdot\vec B}{v}$.

	Les formules de Frenet g\'en\'erales sont
	\begin{align}
		\frac{1}{v}\vec T'&= \kappa\vec N\\
		\frac{1}{v}\vec N'&= -\kappa\vec T+\tau\vec B\\
		\frac{1}{v}\vec B'&= -\tau\vec N
	\end{align}

	\begin{exem}[Tractrice]
		$\vec\alpha(\theta) = \left( \cos\theta+\ln\tan\frac{\theta}{2}, \sin\theta \right)$, avec $\theta \in [\sfrac{\pi}{2}, \pi[$. $\vec\alpha'(\theta) = \left( -\sin\theta+\csc\theta, \cos\theta \right)$. $v(\theta) = \norme{\vec\alpha'(\theta)} = \sqrt{(\sin\theta+\csc\theta)^2+\cos^2\theta} = \sqrt{1-2+\csc^2\theta} = \sqrt{\csc^2\theta-1} = -\cot\theta$.
		\begin{align*}
			\begin{split}
				\vec T(\theta)&= \dfrac{\vec\alpha'(\theta)}{v(\theta)}\\
				&= -\dfrac{\sin\theta}{\cos\theta} \left( -\sin\theta+\csc\theta, \cos\theta \right)\\
				&= \left( \dfrac{\sin^2\theta}{\cos\theta}-\dfrac{1}{\cos\theta}, -\sin\theta \right)\\
				&= \left( -\cos\theta, -\sin\theta \right)
			\end{split}
			&
			\begin{split}
				\kappa\vec N&= \frac{1}{v(\theta)} \vec T'(\theta)\\
				&= -\tan\theta \left( \sin\theta, -\cos\theta \right)\\[.5em]
				\Rightarrow\kappa&= \norme{\kappa\vec N} = -\tan\theta\\
				\vec N&= \frac{\kappa\vec N}{\kappa} = \left( \sin\theta, -\cos\theta \right)
			\end{split}
		\end{align*}
	\end{exem}

	\begin{exem}
		$\vec\alpha(t) = (3t-t^3, 3t^2, 3t+t^3)$. $\vec\alpha'(t) = (3-3t^2, 6t, 3+3t^2)$. $v(t) = \norme{\vec\alpha'(t)} = 3\sqrt{2}(1+t^2)$.
		\begin{align*}
			\vec T(t)&= \frac{\vec\alpha'(t)}{v(t)}\\
			&= \frac{1}{\sqrt{2}}\left( \frac{1-t^2}{1+t^2}, \frac{2t}{1+t^2}, 1 \right)
		\end{align*}
		\begin{align*}
			\kappa\vec N&= \frac{1}{v(t)} \frac{d\vec T}{dt}\\
			&= \frac{1}{3\sqrt{2}(1+t^2)} \left( \cdots \right)\\
			&= \frac{1}{3\sqrt{2}(1+t^2)^2} \left( \frac{-4t}{(1+t^2)^2}, \frac{2(1-t)^2}{(1+t^2)^2}, 0 \right)
		\end{align*}
		\begin{align*}
			\kappa(t)&= \norme{\kappa\vec N}\\
			&= \frac{1}{3(1+t^2)^2}
		\end{align*}
		\begin{align*}
			\vec N(t)&= \frac{\kappa\vec N}{\kappa}\\
			&= \left( -\frac{2t}{(1+t^2)^2}, \frac{1-t^2}{(1+t^2)^2}, 0 \right)
		\end{align*}
		\begin{align*}
			\vec B&= \vec T \times \vec N\\
			&= \frac{1}{\sqrt{2}} \left( -\frac{1-t^2}{1+t^2}, -\frac{2t}{1+t^2}, 1 \right)
		\end{align*}
		\begin{align*}
			-\tau\vec N&= \frac{1}{v}\vec B'\\
			&= -\frac{1}{3(1+t^2)^2} \left( -\frac{2t}{1+t^2}, \frac{1-t^2}{1+t^2}, 0 \right)\\
			\Rightarrow\tau&= \frac{1}{3(1+t^2)^2}\\
			&= \kappa
		\end{align*}
	\end{exem}

	$\vec\alpha'=v\vec T \Longrightarrow \vec\alpha'' = v'\vec T+v\vec T' = v'\vec T+v(v\kappa\vec N) = v'\vec T+v^2\kappa\vec N$.
	\begin{align*}
		\vec\alpha' \times \vec\alpha''&= v\vec T \times (v'\vec T+v^2\kappa\vec N)\\
		&= v^3\kappa\vec T \times \vec N\\
		&= v^3\kappa\vec B\\
		\shortintertext{donc,}
		\norme{\vec\alpha' \times \vec\alpha''}&= v^3\kappa\\
		\kappa&= \frac{\norme{\vec\alpha' \times \vec\alpha''}}{\norme{\vec\alpha'}^3}
	\end{align*}

	\begin{exer}
		\begin{proof}[Montrer que $\vec\alpha$ est lin\'eaire si, et seulement si, $\kappa=0$]
			\begin{align*}
				\kappa(s)=0, \forall s&\Leftrightarrow \norme{\vec T'(s)}=0, \forall s\\
				&\Leftrightarrow \vec T'(s)=0, \forall s\\
				&\Leftrightarrow \vec T(s)=\vec T_0 \text{ const}, \forall s\\
				&\Leftrightarrow \vec\alpha(s)=s\vec T_0+\vec\alpha_0
			\end{align*} o\`u $\vec\alpha_0$ est constante.
		\end{proof}
	\end{exer}

	\begin{exer}
		Soit $\vec\alpha$ une courbe dont toutes les lignes tangentes ($\{\vec\alpha(s)+c\vec T(s) \mid c \in \RR\}$) passent par un m\^eme point, $\vec x_0$, fixe. Que peut-on dire \`a propos de $\vec\alpha$?
		\begin{proof}[D\'emontrons que $\vec\alpha$ est une droite]
			Sans perte de g\'en\'eralit\'e, supposons que $\vec\alpha$ est PPLA et $\vec x_0=0$. On peut donc r\'e\'ecrire l'hypoth\`ese comme $\vec\alpha(s)$ et $\vec T(s)$ sont colin\'eaires, c'est-\`a-dire, $\vec\alpha(s)=\lambda(s)\vec T(s)$, o\`u $\lambda(s) \in \RR$.
			\begin{align*}
				\vec\alpha'(s)&= \lambda'(s)\vec T(s) + \lambda(s)\vec T'(s)\\
				\vec T&= \lambda'\vec T + \lambda\kappa\vec N\\
				(\lambda'-1)\vec T + \lambda\kappa\vec N&= 0
			\end{align*}
			Puisque $\vec T$ et $\vec N$ sont lin\'eairement ind\'ependants, on a que $\lambda'=0$ et $\lambda\kappa=0$. Alors, $\lambda(s)=s+c$ et $\lambda\kappa=0$, donc $\kappa=0$. Ainsi, $\vec\alpha$ est lin\'eaire.
		\end{proof}
	\end{exer}

	\begin{exer}
		\begin{proof}[Montrer que]
			une courbe est planaire (contenue dans un plan $P=\{(x,y,z)\in\RR^3 \mid ax+by+cz=d\}$ pour  $a,b,c,d\in\RR$) si, et seulement si, $\tau(s)=0$, $\forall s$. De plus, toute courbe planaire de courbure constante $\kappa>0$ est une portion de cercle de rayon $\sfrac{1}{\kappa}$.
			\begin{subproof}[$\Rightarrow$]
				Supposons que $\vec\alpha:I\to\RR^3$ est planaire et PPLA. Soit $P=\{(x,y,z)\in\RR^3 \mid ax+by+cz=d\} = \{\vec v\in\RR^3 \mid \vec v \cdot \vec n=d\}$ le plan qui contient $\vec\alpha$, o\`u $\vec n\in\RR^3$ et $\norme{\vec n}=1$.

				On a que $\vec\alpha(s)\in P$, $\forall s$, donc
				\begin{align*}
					\vec\alpha(s) \cdot \vec n&= d\\
					\vec\alpha'(s) \cdot \vec n&= 0\\
					\vec T \cdot \vec n&= 0\\
					\vec T' \cdot \vec n&= 0\\
					\shortintertext{par les formules de Frenet}
					\kappa\vec N \cdot \vec n&= 0
				\end{align*}
				Comme $\vec n$ est orthogonale \`a $\vec T$ et $\vec N$, $\vec n$ est colin\'eaire \`a $\vec B$. Comme $\norme{\vec n}=1$, $\vec n = \pm\vec B$.

				Puisque $\vec n$ est constant, on a que $\vec B'=0$, donc par Frenet, $\vec B'=-\tau\vec N=0$, alors $\tau=0$.
			\end{subproof}
			\begin{subproof}[$\Leftarrow$]
				Supposons que $\tau(s)=0, \forall s$. Alors, $\vec B'=-\tau\vec N=0$, donc $\vec B=B_0$ constant.

				Soit $\vec n=\vec B_0$. Soit $f(s)=\vec\alpha(s) \cdot \vec n$. Alors
				\begin{align*}
					f'(s)&= \vec\alpha'(s) \cdot \vec n\\
					&= \vec T \cdot \vec n\\
					&= \vec T \cdot \vec B_0\\
					&= \vec T \cdot \vec B\\
					&= 0
				\end{align*}
				Donc, $f(s)=d$, pour une constante $d$. Alors, $\vec\alpha(s) \cdot \vec n=d, \forall s$, c'est-\`a-dire, $\vec\alpha(s)\in P=\{\vec v\in\RR^3 \mid \vec v \cdot \vec n=d\}$.
			\end{subproof}

			Supposons que $\vec\alpha$ est planaire et que $\kappa(s)=\kappa_0$, $\forall s$, o\`u $\kappa_0 \in \RR$ est constante.

			Soit $\vec\beta(s) = \vec\alpha(s)+\frac{1}{\kappa_0}\vec N(s)$.
			\begin{align*}
				\vec\beta'(s)&= \vec\alpha'(s) + \frac{1}{\kappa_0} \vec N'(s)\\
				&= \vec T + \frac{1}{\kappa_0} \left( -\kappa\vec T + \tau\vec B \right)\\
				&= \vec T - \vec T\\
				&= 0
			\end{align*}
			Alors, $\vec\beta(s)=\vec P_0$, $\forall s$, o\`u $\vec P_0\in\RR^3$ est constant.
			\begin{align*}
				\norme{\vec\alpha(s) - \vec P_0}&= \norme{\vec\alpha(s) - \vec\beta(s)}\\
				&= \norme{-\frac{1}{\kappa_0} \vec N(s)}\\
				&= \frac{1}{\kappa_0}
			\end{align*}
			Ainsi, $\vec\alpha(s) \subseteq \left\{ \vec v\in\RR^3  ~\middle|~ \begin{array}{l}
				\vec v \cdot \vec n=d\\
				\norme{\vec v-\vec P_0}=\frac{1}{\kappa_0}
			\end{array}\right\}$, c'est-\`a-dire, $\vec\alpha(s)$ est dans le cercle de rayon $\sfrac{1}{\kappa_0}$ centr\'e en $\vec P_0$ dans le plan $P$.
		\end{proof}
	\end{exer}

	\begin{exem}[H\'elice]
		$\vec\alpha(s) = \left( a\cos\frac{s}{c}, a\sin\frac{s}{c}, \frac{bs}{c} \right)$, avec $a,b\in\RR$. $\kappa=\frac{a}{c^2}$. $\tau=\frac{b}{c^2}$.

		Soit $\vec A=(0,0,1)\in\RR^3$. $\vec T \cdot \vec A = \left( -\frac{a}{c}\sin\frac{s}{c}, \frac{a}{c}\cos\frac{s}{c}, \frac{b}{c} \right) \cdot (0,0,1) = \frac{b}{c}$, constant.

		Une \emph{h\'elice g\'en\'eralis\'ee} est une courbe $\vec\alpha$ avec $\kappa\neq0$ o\`u $\vec T$ fait un angle constant avec une direction fixe. C'est-\`a-dire, $\exists\vec A\in\RR^3$ unitaire t.q. $\frac{\vec T \cdot \vec A}{\norme{\vec T} \norme{\vec A}}=\cos\theta$, o\`u $\theta$ est constant. C'est-\`a-dire, $\vec T \cdot \vec A$ est constant.

		Soit $P = \left\{ \vec v\in\RR^3 ~\middle|~ \vec v \cdot \vec A=0 \right\}$ le plan d\'etermin\'e par $\vec A$.

		Soient $\vec X, \vec Y \in P$ une base orthonorm\'ee. Donc $\vec X, \vec Y, \vec A$ est une base orthonorm\'ee de $\RR^3$. Alors, $\vec\alpha(s) = u(s)\vec X+v(s)\vec Y+w(s)\vec A$ pour des fonctions $u,v,w$. $\vec T(s) = u'(s)\vec X+v'(s)\vec Y+w'(s)\vec A$. $\vec T(s) \cdot \vec A=w'(s)$ constant, alors $w(s)=a+bs$ pour $a,b\in\RR$. $\vec\alpha(s) = \underbrace{u(s)\vec X+v(s)\vec Y}_{\text{courbe planaire dans $P$}} + \underbrace{(a+bs)\vec A}_{\text{partie lin\'eaire}}$.
	\end{exem}

	\begin{figure}[h]
		\centering
		\begin{tikzpicture}[samples=50, domain=-3:3, variable=\t]
			\begin{scope}[canvas is xz plane at y=0]
				\fill[lightgray, fill opacity=.3] (-4,-2) rectangle (4,2);
			\end{scope}
			\draw[thick, red] plot (\t,0,{sin(\t r)});
			\foreach \x in {-3,-2,...,3}{
				\draw[help lines] (\x,0,{sin(\x r)}) -- (\x,3,{sin(\x r)});
			}
			\foreach \y in {0,1,2,3}{
				\draw[help lines] plot (\t,\y,{sin(\t r)});
			}
			\draw[thick, blue] plot (\t,{(\t+3)/2},{sin(\t r)});
			\node[blue] () at (3.2,3.1,{sin(3.2 r)}) {$\vec\alpha$};
			\draw[-latex, thick] (-3,0,{sin(-3 r)}) -- ++(0,1,0) node[above left] {$\vec A$};
			\draw (-3,.4,{sin(-3 r)}) -- (-2.8,.1,{sin(-2.8 r)});
			\node () at (-2.8,.4,{sin(-2.8 r)}) {$\theta$};
		\end{tikzpicture}
		\caption{Exemple d'une h\'elice g\'en\'eralis\'ee}
	\end{figure}

	\begin{prop}
		Une courbe $\vec\alpha$ avec $\kappa\neq0$ est une h\'elice g\'en\'eralis\'ee si, et seulement si, $\frac{\tau}{\kappa}$ est constant.
		\begin{proof}
			\begin{subproof}[$\Rightarrow$:]
				Soit $\vec\alpha$ une h\'elice g\'en\'eralis\'ee PPLA. Soit $\vec A\in\RR^3$ un vecteur unitaire t.q. $\vec T \cdot \vec A=\cos\theta$ pour $\theta$ constant.
				\begin{align*}
					\vec T' \cdot \vec A&= 0\\
					\shortintertext{par les formules de Frenet,}
					\kappa\vec N \cdot \vec A&= 0\\
					\shortintertext{comme $\kappa\neq0$,}
					\vec N \cdot \vec A&= 0\tag{1}\\
					\vec N' \cdot \vec A&= 0\\
					\shortintertext{par les formules de Frenet,}
					(-\kappa\vec T+\tau\vec B) \cdot \vec A&= 0\tag{2}
				\end{align*}

				De (1), on a $\vec A$ est dans le plan g\'en\'er\'e par $\vec T$ et $\vec B$. Comme l'angle entre $\vec T$ et $\vec A$ est $\theta$, $\vec A = \cos\theta\vec T \pm \sin\theta\vec B$.

				De (2), on a $(-\kappa\vec T+\tau\vec B) \cdot (\cos\theta\vec T \pm \sin\theta\vec B)=0$, alors $-\kappa\cos\theta\pm\tau\sin\theta=0$, donc $\frac{\tau}{\kappa}=\pm\cot\theta$, qui est constant.
			\end{subproof}
			\begin{subproof}[$\Leftarrow$:]
				Supposons que $\frac{\tau}{\kappa}$ est constant. Soit $\theta\in\RR$ t.q. $\frac{\tau}{\kappa}=\cot\theta$. Soit $\vec A=\cos\theta\vec T\pm\sin\theta\vec B$. On a que $\vec A \cdot \vec T=\cos\theta$. Il suffit de montrer que $\vec A$ est constant.
				\begin{align*}
					\vec A'&= \cos\theta\vec T' + \sin\theta\vec B'\\
					\shortintertext{par les formules de Frenet,}
					&= \cos\theta\kappa\vec N + \sin\theta(-\tau\vec N)\\
					&= \left( \cos\theta\kappa-\sin\theta\tau \right) \vec N\\
					&= \kappa\sin\theta \left( \cot\theta-\frac{\tau}{\kappa} \right) \vec N\\
					&= 0
				\end{align*}
			\end{subproof}
		\end{proof}
	\end{prop}

	\begin{exem}
		$\vec\alpha(t) = \left( 3t-t^3, 3t^2, 3t+t^3 \right)$. $\kappa(t) = \tau(t) = \frac{1}{3(1+t^2)^2}$. Alors, $\frac{\tau}{\kappa} = 1 = \cot\frac{\pi}{4}$ est constant. $\theta=\frac{\pi}{4}$, alors $\vec A=\cos\theta\vec T+\sin\theta\vec B = \frac{1}{\sqrt{2}} \frac{1}{\sqrt{2}}\left( \frac{1-t^2}{1+t^2}, \frac{2t}{1+t^2}, 1 \right) + \frac{1}{\sqrt{2}} \frac{1}{\sqrt{2}} \left( -\frac{1-t^2}{1+t^2}, \frac{-2t}{1+t^2}, 1 \right) = (0,0,1)$. $\vec\alpha$ fait un angle de $\frac{\pi}{4}$ avec $(0,0,1)$.
	\end{exem}

	\section{Forme locale canonique}
	Soit $\vec\alpha$ une courbe PPLA. Sans perte de g\'en\'eralit\'e, $\vec\alpha(0)=\vec0\in\RR^3$.

	Par la formule de Taylor,
	\begin{align*}
		\vec\alpha(s)&= \vec\alpha(0) + s\vec\alpha'(0) + \frac{s^2}{2}\vec\alpha''(0) + \frac{s^3}{6}\vec\alpha'''(0) + \cdots\\
		&= s\vec T(0) + \frac{s^2}{2}\vec T'(0) + \frac{s^3}{6}\vec T''(0) + \cdots\\
		&= s\vec T_0 + \frac{s^2}{2}\kappa(0)\vec N(0) + \frac{s^3}{6} \left(\kappa\vec N \right)'(0) + \cdots\\
		&= s\vec T_0 + \frac{s^2}{2}\kappa_0\vec N_0 + \frac{s^3}{6} \left(\kappa'_0\vec N_0 + \kappa_0\vec N'_0 \right) + \cdots\\
		&= s\vec T_0 + \frac{s^2}{2} \kappa_0\vec N_0 + \frac{s^3}{6} \left( \kappa'_0\vec N_0 + \kappa_0(-\kappa_0\vec T_0+\tau_0\vec B_0) \right) + \cdots\\
		&= \vec T_0 \left( s-\frac{\kappa_0^2}{6}s^3+\cdots \right) + \vec N_0 \left( \frac{\kappa_0}{2}s^2 + \frac{\kappa'_0}{6}s^3+\cdots \right) + \vec B_0 \left( \frac{\kappa_0\tau_0}{6}s^3+\cdots \right)
	\end{align*}

	\begin{dfn}[Plan osculateur]
		Le \emph{plan osculateur} est le plan g\'en\'er\'e par $\vec T$ et $\vec N$.

		La projection de $\vec\alpha$ sur le plan osculateur est $\left( s-\frac{\kappa_0^2}{6}s^3+\cdots, \frac{\kappa_0}{2}s^2+\frac{\kappa'_0}{6}s^3+\cdots \right)$.
	\end{dfn}

	On cherche $r$ t.q. le cercle de rayon $r$ passant par $(0,0)$ \`a $s=0$ centr\'e en $(0,r)$ soit la meilleure approximation de $\vec\alpha$. $(0,r)+(r\sin t, -r\cos t) = (r\sin t, r-r\cos t)$. En reparam\'etrisant par longueur d'arc, $\vec\beta(s) = \left( r \sin\frac{s}{r}, r\left( 1-\cos\frac{s}{r} \right) \right)$. Soit $\vec\beta(s) = r\sin\frac{s}{r}\vec T_0 + r\left( 1-\cos\frac{s}{r} \right)\vec N_0$, $\vec\beta(0)=0=\vec\alpha(0)$.

	Par Taylor,
	\begin{align*}
	\vec\beta(s)&= r\left( \frac{s}{r} - \frac{s^3}{6r^3} - \cdots \right)\vec T_0 + r\left( 1-\left( 1 - \frac{s^2}{2r^2} + \cdots \right) \right)\vec N_0\\
	&= \vec T_0\left( s-\frac{\sfrac{1}{r^2}}{6}s^3+\cdots \right) + \vec N_0 \left( \frac{\sfrac{1}{r}}{2}s^2+\cdots \right)
	\end{align*}

	$r=\sfrac{1}{\kappa_0}$ est le rayon pour lequel $\vec\beta(s)$ approxime le mieux la courbe $\vec\alpha(s)$ pr\`es de $s=0$.

	Donc $\kappa(0)=\sfrac{1}{r}$, o\`u $r$ est le rayon du cercle qui approxime le mieux $\vec\alpha$ pr\`es de $s=0$.

	$\vec\beta(s)$ est l'unique cercle t.q.
	\begin{align*}
	\vec\beta(0)&= \vec\alpha(0)\\
	\vec\beta'(0)&= \vec\alpha'(0)\\
	\vec\beta''(0)&= \vec\alpha''(0)
	\end{align*}

	On peut interpr\'eter ce r\'esultat comme $\kappa(s)$ est l'inverse du rayon de l'unique cercle $\vec\beta$ dans $\RR^3$ t.q.
	\begin{align*}
	\vec\beta(s)&= \vec\alpha(s)\\
	\vec\beta'(s)&= \vec\alpha'(s)\\
	\vec\beta''(s)&= \vec\alpha''(s)
	\end{align*}

	\begin{dfn}
		Une \emph{transformation rigide} (ou euclidienne) est la composition d'une rotation et d'une translation. \[ \Psi : \RR^3\to\RR^3, \Psi(\vec x) = A\vec x+\vec b \] o\`u $A$ est une matrice de rotation, c'est-\`a-dire, $A^TA=I$ et $\det A=1$, et $\vec b\in\RR^3$ est constant.

		Deux courbes PPLA $\vec\alpha$ et $\vec\alpha^*$ sont \emph{congruentes} si $\exists \Psi$ une transformation rigide t.q. $\vec\alpha^*(s)=\Psi(\vec\alpha(s))$, $\forall s$.
	\end{dfn}

	\begin{thm}[Th\'eor\`eme fondamental des courbes]~

		\begin{enumerate}
			\item Soient $\vec\alpha$ et $\vec\alpha^*$ deux courbes PPLA avec $\kappa\neq0$ et $\kappa^*\neq0$. Alors, $\vec\alpha$ et $\vec\alpha^*$ sont congruentes si, et seulement si, $\kappa=\kappa^*$ et $\tau=\tau^*$;
			\item Soit $\kappa,\tau:[0,L]\to\RR$ des fonctions continues avec $\kappa>0$. Alors, il existe une courbe $\vec\alpha$ PPLA avec $\kappa$ comme courbure et $\tau$ comme torsion.
		\end{enumerate}
		\begin{subproof}[Id\'ee de la d\'emonstration]
			\begin{align*}
				\begin{split}
					\vec T'&= \kappa\vec N\\
					\vec N'&= -\kappa\vec T+\tau\vec B\\
					\vec B'&= -\tau\vec N
				\end{split}\tag{$*$}
			\end{align*}

			Soit $\vec X = \begin{pmatrix}
				T_1&T_2&T_3&N_1&N_2&N_3&B_1&B_2&B_3
			\end{pmatrix}^T$, o\`u $\vec T=\begin{pmatrix}
				T_1&T_2&T_3
			\end{pmatrix}$, \dots. $\vec X'=A\vec X$, o\`u \[ A = \left[\begin{array}{ccc|ccc|ccc}
				&&&\kappa&&&&&\\
				&&&&\kappa&&&&\\
				&&&&&\kappa\\\hline
				-\kappa&&&&&&\tau\\
				&-\kappa&&&&&&\tau\\
				&&-\kappa&&&&&&\tau\\\hline
				&&&-\tau&&&&&\\
				&&&&-\tau&&&&\\
				&&&&&-\tau&&&
			\end{array}\right] \in \RR^{9\times9} \]
			un syst\`eme d'\'equations diff\'erentielles du premier ordre. Par le th\'eor\`eme fondamental d'existence et d'unicit\'e, les \'equations $(*)$ ont une unique solution pour chaque condition initiale $\vec T_0, \vec N_0, \vec B_0$.

			Choisir une base orthonorm\'ee $\vec T_0, \vec N_0, \vec B_0$. Soit $\vec T, \vec N, \vec B$ la solution. Soit $\vec\alpha(s)=\int_{0}^{s}\vec T(u)du$.
		\end{subproof}
	\end{thm}

	\begin{exem}
		Trouver une courbe avec $\kappa(s) = \frac{1}{1+s}$ et $\tau(s)=0$.
		\begin{align*}
			\vec T'&= \frac{1}{1+s}\vec N\\
			\vec N'&= -\frac{1}{1+s}\vec T\\
			\vec B'&= 0
		\end{align*}

		Soit $u=\ln(1+s)$, c'est-\`a-dire, $1+s=e^u$.
		\begin{align*}
			\begin{split}
				\frac{d\vec T}{du}&= \frac{\sfrac{d\vec T}{ds}}{\sfrac{du}{ds}}\\
				&= \frac{\frac{1}{1+s}\vec N}{\frac{1}{1+s}}\\
				&= \vec N
			\end{split}\\
			\begin{split}
				\frac{d\vec N}{du}&= \frac{\sfrac{d\vec N}{ds}}{\sfrac{du}{ds}}\\
				&= \frac{-\frac{1}{1+s}\vec T}{\frac{1}{1+s}}\\
				&= -\vec T
			\end{split}\\
			\begin{split}
				\frac{d\vec B}{du}&= 0
			\end{split}
		\end{align*}

		Alors,
		\begin{align*}
			\vec T&= \cos u\vec c_1+\sin u\vec c_2\\
			\vec N&= -\sin u\vec c_1+\cos u\vec c_2\\
			\vec B&= \vec c_3
		\end{align*}
		o\`u $\vec c_1, \vec c_2, \vec c_3 \in \RR^3$ orthonorm\'es.

		Posons $\vec c_1=(1,0,0)$, $\vec c_2=(0,1,0)$, $\vec c_3=(0,0,1)$. $\vec T(u)=(\cos u,\sin u,0) = (\cos(\ln(1+s)), \sin(\ln(1+s)), 0)$.

		\begin{align*}
			\vec\alpha(s)&= \int_{0}^{s} \vec T(x)dx\\
			&= \int_{0}^{s} \left( \cos(\ln(1+x)), \sin(\ln(1+x)), 0 \right) dx\\
			&= \left( \int_{0}^{s}\cos(\ln(1+x))dx, \int_{0}^{s}\sin(\ln(1+x))dx, 0 \right)\\
			&= \left( \frac{1}{2}(-1+(1+s)(\cos\ln(1+s))+\sin(\ln(1+s))), \frac{1}{2}(1-(1+s)(\cos(\ln(1+s))-\sin(\ln(1+s)))), 0 \right)
		\end{align*}
	\end{exem}

	\section{Courbes planaires}
	Soit $\vec\alpha:[0,L]\to\RR^2$ une courbe planaire PPLA.
	\begin{align*}
		\vec\alpha(s)&= (x(s), y(s))\\
		\vec T(s)&= (x'(s), y'(s))
	\end{align*}
	$\norme{\vec T(s)}=1$, donc $\vec T(s)$ est un point sur le cercle unitaire.

	\begin{dfn}
		$\vec N(s)$ est l'unique vecteur unitaire dans $\RR^2$ t.q. $\vec N(s) \cdot \vec T(s)$ et $\{\vec T(s), \vec N(s)\}$ est orient\'ee positivement, c'est-\`a-dire,
		\begin{itemize}
			\item $\vec T \times \vec N=(0,0,1)$;
			\item $\vec N(s)$ est la rotation de $\vec T(s)$ de 90\degree\ dans le sens trigonom\'etrique;
			\item si $\vec T=(a,b)$, alors $\vec N=(-b,a)$.
		\end{itemize}
	\end{dfn}

	\begin{rmq}
		On n'a pas besoin de supposer que $\vec T'\neq0$ pour d\'efinir $\vec N$ dans ce cas.
	\end{rmq}

	\begin{dfn}
		$\kappa=\vec T' \cdot \vec N$.
	\end{dfn}

	\begin{rmq}
		$\vec T'=\kappa\vec N$, car
		\begin{align*}
			\vec T'&= (\vec T' \cdot \vec T)\vec T + (\vec T' \cdot \vec N)\vec N\\
			\shortintertext{comme $\norme{\vec T}=1$ et $\vec T' \cdot \vec N=\kappa$,}
			&= \kappa\vec N
		\end{align*}
	\end{rmq}

	$\kappa$ peut \^etre $>0$, $=0$ ou $<0$. Si $\kappa>0$, alors $\vec T$ tourne dans le sens trigonom\'etrique. Si $\kappa<0$, alors $\vec T$ tourne dans le sens horaire. Si $\kappa=0$, alors $\vec T$ ne tourne pas, la courbe peut \^etre une droite ou il s'agit d'un point critique.

	\begin{dfn}
		On a que $\vec T(s) = (\cos\theta(s), \sin\theta(s))$, o\`u $\theta:[0,L]\to\RR$ est continue.

		$\vec T' = (-\theta'\sin\theta, \theta'\cos\theta) = \theta'(-\sin\theta, \cos\theta) = \theta'\vec N$.

		Alors, $\kappa = \vec T' \cdot \vec N = \theta'\vec N \cdot \vec N = \theta'$.
		\[ \kappa(s) = \theta'(s) \quad\Longrightarrow\quad \int_{0}^{L}\kappa(s) ds = \theta(L)-\theta(0)\]
	\end{dfn}
\end{document}
